\documentclass[11pt,a4paper]{article}

% ----- Pacotes -----
\usepackage[utf8]{inputenc}
\usepackage[T1]{fontenc}
\usepackage[brazil]{babel}
\usepackage[margin=2.5cm]{geometry}
\usepackage[table]{xcolor}        % cor + listras de linha em tabelas
\usepackage{booktabs}
\usepackage{tabularx}             % colunas elásticas (texto que quebra)
\usepackage{ragged2e}            % alinhamento em colunas X
\usepackage{amsmath}
\usepackage{array}
\usepackage{enumitem}
\usepackage[font=small,labelfont=bf,skip=5pt]{caption}
\usepackage[nobottomtitles*]{titlesec}   % evita títulos no rodapé da página

% ----- Cores (definidas antes de quem as usa) -----
\definecolor{azulsec}{RGB}{31,59,87}     % azul dos títulos
\definecolor{cinzahdr}{RGB}{31,59,87}    % cabeçalho de tabela
\definecolor{cinzalin}{RGB}{238,242,246} % listra clara das linhas

\usepackage{hyperref}
\hypersetup{colorlinks=true,linkcolor=azulsec,urlcolor=azulsec,citecolor=azulsec}

\setlength{\parskip}{0.5em}
\setlength{\parindent}{0pt}
\renewcommand{\arraystretch}{1.25}
\setlength{\tabcolsep}{6pt}

% Evita linhas órfãs/viúvas e títulos colados no fim da página
\widowpenalty=10000
\clubpenalty=10000
\setlength{\headsep}{12pt}

% Títulos de seção em azul, com espaçamento que segura o texto seguinte
\titleformat{\section}{\Large\bfseries\color{azulsec}}{\thesection}{0.6em}{}
\titleformat{\subsection}{\large\bfseries\color{azulsec}}{\thesubsection}{0.6em}{}
\titlespacing*{\section}{0pt}{2.2ex plus 1ex minus .2ex}{1.2ex plus .2ex}
\titlespacing*{\subsection}{0pt}{1.8ex plus 1ex minus .2ex}{0.9ex plus .2ex}

% Estilo padrão de cabeçalho de tabela (texto branco sobre azul)
\newcommand{\hdr}[1]{\textcolor{white}{\textbf{#1}}}
\newcolumntype{Y}{>{\RaggedRight\arraybackslash}X}  % coluna X justificada à esquerda

% ----- Capa -----
\title{\textbf{Mortalidade por Causas Evitáveis no Brasil}\\[0.3em]
\large Desigualdade socioeconômica, padrões municipais e perfis de causa\\(2022--2024)}
\author{Projeto Data Major --- Tópicos em Banco de Dados\\IESB}
\date{Maio de 2026}

\begin{document}

\maketitle
\thispagestyle{empty}
\clearpage

\begin{abstract}
\noindent Este trabalho investiga a relação entre a mortalidade por causas evitáveis e as
condições socioeconômicas dos municípios brasileiros. A partir de um pipeline de dados
completo (ETL), integramos os microdados de óbitos do SIM/DATASUS (2022--2024) a indicadores
de desenvolvimento humano, renda, desigualdade e investimento em saúde. Aplicamos a Lista
Brasileira de Causas Evitáveis (LBCE) para identificar $\approx$1,37 milhão de óbitos
evitáveis na faixa de 5 a 74 anos. Sobre esse conjunto realizamos três análises
complementares: (i) regressão multivariada (OLS, Random Forest e XGBoost) para medir a
associação entre variáveis socioeconômicas e a taxa de mortalidade evitável;
(ii) clusterização (\emph{K-Means}) dos municípios em quatro perfis socioeconômicos; e
(iii) mineração das causas básicas (CID-10) para caracterizar \emph{do que} morre cada perfil.
Para garantir comparabilidade entre municípios com estruturas etárias distintas, calculamos a
taxa de mortalidade \textbf{padronizada por idade} (padronização direta com a população do
Censo 2022), adotada como desfecho principal. Os resultados confirmam associação
estatisticamente significativa entre desigualdade (Gini) e mortalidade evitável, e revelam uma
assinatura socioeconômica nas causas: municípios mais pobres concentram doenças infecciosas e
negligenciadas, enquanto os mais desenvolvidos concentram doenças crônicas. A padronização por
idade corrige um paradoxo das taxas brutas (Sul/Sudeste apareciam com mortalidade maior por
terem população mais envelhecida) e alinha os sinais da regressão à hipótese socioeconômica.
\end{abstract}

\clearpage
\tableofcontents
\clearpage

% ===================================================================
\section{Introdução}

\textbf{Problema.} Causas evitáveis de morte são aquelas que poderiam ser prevenidas ou
tratadas por ações do Sistema Único de Saúde (SUS): imunização, diagnóstico precoce, atenção
básica e políticas de saúde pública. A persistência de altas taxas dessas mortes no Brasil é
considerada um indicador de iniquidade --- pessoas morrem de causas que o sistema de saúde,
em tese, sabe prevenir.

\textbf{Hipótese.} A taxa de mortalidade por causas evitáveis de um município está associada
ao seu nível de desenvolvimento socioeconômico (renda, IDH-M, desigualdade, saneamento,
escolaridade) e ao esforço público em saúde.

\textbf{Objetivo.} Construir um pipeline de dados reprodutível e aplicar técnicas de
modelagem e mineração para: (1) quantificar a associação entre condições socioeconômicas e a
mortalidade evitável; (2) identificar perfis típicos de município; e (3) caracterizar o
padrão de causas de cada perfil.

% ===================================================================
\section{Fontes de Dados}

A Tabela~\ref{tab:fontes} resume as cinco fontes integradas, todas públicas.

\begin{table}[htbp]
\centering\small
\caption{Fontes de dados utilizadas.}
\label{tab:fontes}
\rowcolors{2}{cinzalin}{white}
\begin{tabularx}{\textwidth}{l Y >{\RaggedRight\arraybackslash}p{4.0cm}}
\rowcolor{cinzahdr}
\hdr{Fonte} & \hdr{Conteúdo} & \hdr{Acesso} \\
\midrule
SIM / DATASUS & Microdados de óbitos (CID-10), por UF e ano & FTP (arquivos \texttt{.dbc}) \\
IBGE / Sidra  & População municipal (Censo 2022; estimativas) & API (tab.\ 4714 e 6579) \\
IBGE / Sidra  & População por faixa etária $\times$ município & API (tab.\ 9514) \\
IBGE / Sidra  & PIB municipal & API (tab.\ 5938) \\
SICONFI / Tesouro & Despesa municipal em saúde (função 10) & API DCA-Anexo I-E \\
IPEA / Atlas DH & IDH-M, Gini, renda, pobreza, saneamento & API (séries \texttt{ADH\_*}) \\
\bottomrule
\end{tabularx}
\end{table}

\textbf{Tratamento de lacunas temporais.} O painel cobre os anos de 2022, 2023 e 2024, mas
nem toda fonte tem cobertura completa nesses anos. Adotamos substituições explícitas,
registradas em uma coluna \texttt{ANO\_FONTE} para auditoria:
\begin{itemize}[nosep]
  \item \textbf{População 2023}: o IBGE não publicou estimativa; usamos o Censo 2022 como proxy.
  \item \textbf{PIB 2024}: indisponível por defasagem contábil ($\approx$2 anos); usamos o PIB de 2023.
  \item \textbf{Atlas DH}: calculado apenas com base no Censo 2010; é tratado como
        \emph{invariante no tempo} (repetido nos três anos).
\end{itemize}

% ===================================================================
\section{Metodologia}

\subsection{Classificação das causas evitáveis (LBCE)}

Utilizamos a \textbf{Lista Brasileira de Causas de Mortes Evitáveis} de 5 a 74 anos
(Malta et al., 2007; atualização 2010), que organiza códigos CID-10 segundo o tipo de
intervenção do SUS capaz de evitar a morte. A lista oficial possui cinco grupos; neste
trabalho utilizamos \textbf{quatro}, descritos na Tabela~\ref{tab:lbce}.

\begin{table}[htbp]
\centering\small
\caption{Grupos da LBCE utilizados na análise.}
\label{tab:lbce}
\rowcolors{2}{cinzalin}{white}
\begin{tabularx}{\textwidth}{l Y}
\rowcolor{cinzahdr}
\hdr{Grupo} & \hdr{Descrição} \\
\midrule
\texttt{imunoprevencao}             & Reduzíveis por vacinação \\
\texttt{doencas\_infecciosas}       & Tuberculose, HIV, hepatites, infecções respiratórias, doenças tropicais \\
\texttt{doencas\_nao\_transmissiveis} & Neoplasias, doenças cardiovasculares, diabetes, DPOC \\
\texttt{morte\_materna}             & Gravidez, parto e puerpério \\
\bottomrule
\end{tabularx}
\end{table}

\textbf{Exclusão das causas externas.} O quinto grupo da lista oficial --- causas externas
(acidentes, quedas, agressões, suicídios) --- foi \textbf{deliberadamente excluído}. As
variáveis socioeconômicas disponíveis explicam mal essa mortalidade: acidentes de trânsito,
homicídios e quedas respondem a fatores como segurança pública, infraestrutura viária e
violência interpessoal, ausentes do nosso conjunto de \emph{features}. Incluí-las
introduziria ruído estrutural sem ganho explicativo. O escopo final é, portanto, o das
causas evitáveis \emph{sensíveis ao SUS e ao desenvolvimento socioeconômico}.

\subsection{Pipeline de dados}

O projeto é organizado em notebooks sequenciais (Tabela~\ref{tab:notebooks}), do ETL à
mineração. Cada etapa é idempotente e persiste seus resultados em disco.

\begin{table}[htbp]
\centering\small
\caption{Notebooks do pipeline.}
\label{tab:notebooks}
\rowcolors{2}{cinzalin}{white}
\begin{tabularx}{\textwidth}{l Y}
\rowcolor{cinzahdr}
\hdr{Notebook} & \hdr{Função} \\
\midrule
\texttt{01\_extracao}                 & Download e cache das cinco fontes (com \emph{fallbacks}). \\
\texttt{02\_transformacao}            & Aplicação da LBCE, agregação e construção da \emph{feature matrix} (município $\times$ ano). \\
\texttt{02b\_transformacao\_microdados} & Dataset individualizado (1 óbito por linha) com estatística descritiva. \\
\texttt{03\_regressao}                & Regressão multivariada em painel (OLS, RF, XGBoost). \\
\texttt{04\_clusterizacao}            & Clusterização \emph{K-Means} dos municípios. \\
\texttt{05\_mineracao}                & Caracterização das causas CID-10 por cluster. \\
\texttt{06\_taxa\_padronizada\_idade}  & Padronização etária da taxa e comparação com a bruta. \\
\bottomrule
\end{tabularx}
\end{table}

\textbf{Unidade de análise e \emph{join}.} O SIM identifica o município de residência por um
código de 6 dígitos (\texttt{CODMUNRES}); o IBGE usa 7 dígitos (com dígito verificador). A
junção entre fontes é feita pelos 6 primeiros dígitos. A \emph{feature matrix} final tem uma
linha por município e ano (16.711 linhas; 5.570 municípios $\times$ 3 anos).

\textbf{Variável-resposta.} A taxa bruta de mortalidade evitável por 100 mil habitantes é
\[
\text{taxa bruta} = \frac{\text{óbitos evitáveis (5--74 anos)}}{\text{população}} \times 100{.}000 .
\]
Como desfecho principal, porém, usamos a \textbf{taxa padronizada por idade} (subseção
seguinte), que neutraliza diferenças de estrutura etária entre municípios.

\subsection{Padronização por idade}

A taxa bruta não corrige a estrutura etária: como as causas evitáveis (5--74 anos) são
dominadas por doenças crônicas que matam perto dos 70, municípios mais envelhecidos têm taxa
bruta maior \emph{só por isso}. Calculamos a taxa \textbf{padronizada por idade} (padronização
direta): para cada município, a taxa específica de cada faixa quinquenal é ponderada pelos
pesos da estrutura etária nacional (Censo 2022):
\[
\text{taxa padronizada}_m = \sum_{a} \frac{\text{óbitos}_{m,a}}{\text{população}_{m,a}}\; w_a \times 100{.}000,
\qquad w_a = \frac{\text{pop. nacional da faixa } a}{\text{pop. nacional 5--74}} .
\]
O numerador por faixa vem dos microdados (idade de cada óbito); o denominador por faixa é
extraído do Censo 2022 (IBGE, tabela 9514). É a taxa que cada município teria se tivesse a
estrutura etária do Brasil --- tornando as comparações entre eles justas.

\subsection{Modelagem (regressão em painel)}

Tratamos os três anos como um \emph{painel agrupado} (\emph{pooled}), com o ano incluído como
variável \emph{dummy} para capturar choques comuns de período. Optou-se por não usar efeitos
fixos de município porque a maioria das \emph{features} (Atlas 2010) é invariante no tempo e
seria absorvida por esses efeitos. Três modelos foram comparados:
\begin{itemize}[nosep]
  \item \textbf{OLS} (\emph{statsmodels}), com erros-padrão \emph{cluster-robust} por
        município --- corrige a correlação entre as três observações de cada cidade.
  \item \textbf{Random Forest} e \textbf{XGBoost}, para capturar não-linearidades e interações.
\end{itemize}
A divisão treino/teste usa \texttt{GroupShuffleSplit} por município (80/20), impedindo que o
mesmo município apareça em treino e teste --- o que inflaria artificialmente o $R^2$.

\subsection{Clusterização}

Para a tipologia de municípios, agregamos as três observações de cada município em uma única
linha (média das variáveis temporais; valor único das do Atlas). Aplicamos \emph{K-Means}
sobre sete \emph{features} socioeconômicas padronizadas (IDH-M, Gini, log do PIB \emph{per
capita}, log da despesa em saúde \emph{per capita}, analfabetismo, água encanada e esperança
de vida). O número de clusters ($k=4$) foi escolhido pela combinação de método do cotovelo e
\emph{silhouette}, priorizando a interpretabilidade. A taxa de mortalidade \textbf{não}
entrou nas \emph{features} (evita circularidade): ela é usada apenas para caracterizar os
clusters depois de formados.

\subsection{Mineração das causas}

Cada óbito individual foi associado ao cluster de seu município. As descrições oficiais dos
códigos CID-10 vieram das tabelas auxiliares do DATASUS (\texttt{CID10.DBF}). Para identificar
causas \emph{características} de cada cluster --- e não apenas as mais frequentes --- usamos o
\emph{lift}:
\[
\text{lift}(causa, cluster) = \frac{P(causa \mid cluster)}{P(causa \mid \text{Brasil})},
\]
em que valores acima de 1 indicam sobre-representação no cluster.

% ===================================================================
\section{Resultados}

\subsection{Panorama descritivo}

Dos 4,54 milhões de óbitos registrados no triênio, 2,49 milhões estão na faixa de 5 a 74
anos, dos quais \textbf{1,37 milhão} foram classificados como evitáveis pela LBCE. A
distribuição entre grupos (Tabela~\ref{tab:grupos}) é fortemente dominada pelas doenças não
transmissíveis.

\begin{table}[htbp]
\centering\small
\caption{Óbitos evitáveis por grupo LBCE (triênio 2022--2024).}
\label{tab:grupos}
\rowcolors{2}{cinzalin}{white}
\begin{tabular}{l r r}
\rowcolor{cinzahdr}
\hdr{Grupo} & \hdr{Óbitos} & \hdr{\%} \\
\midrule
Doenças não transmissíveis & 1.109.236 & 81,2 \\
Doenças infecciosas        &   250.735 & 18,4 \\
Morte materna              &     4.827 &  0,4 \\
Imunoprevenção             &     1.669 &  0,1 \\
\midrule
\textbf{Total}             & \textbf{1.366.467} & \textbf{100,0} \\
\bottomrule
\end{tabular}
\end{table}

A taxa bruta média municipal foi de aproximadamente \textbf{226 óbitos evitáveis por 100 mil
habitantes}, com grande dispersão (de $<$50 a $>$650 por 100 mil).

\textbf{Efeito da padronização por idade.} A Tabela~\ref{tab:gradiente} contrasta a taxa
bruta com a padronizada por região. Pela bruta, o gradiente é \emph{invertido} em relação à
hipótese (Sul/Sudeste no topo). A padronização corrige o paradoxo: o \textbf{Nordeste passa a
superar} Sudeste e Sul, e o Norte sobe de 159 para 228 --- coerente com a tese de que regiões
menos desenvolvidas têm maior mortalidade evitável quando se neutraliza a estrutura etária.

\begin{table}[htbp]
\centering\small
\caption{Taxa bruta vs.\ padronizada por idade, por região (média dos municípios, /100k).}
\label{tab:gradiente}
\rowcolors{2}{cinzalin}{white}
\begin{tabular}{l r r}
\rowcolor{cinzahdr}
\hdr{Região} & \hdr{Taxa bruta} & \hdr{Taxa padronizada} \\
\midrule
Norte        & 158,9 & 227,6 \\
Nordeste     & 204,1 & \textbf{247,3} \\
Centro-Oeste & 229,8 & 257,4 \\
Sudeste      & 244,6 & 244,5 \\
Sul          & 256,5 & 241,9 \\
\bottomrule
\end{tabular}
\end{table}

\subsection{Regressão}

Os modelos foram estimados sobre a \textbf{taxa padronizada} (desfecho principal). O
desempenho preditivo é baixo (Tabela~\ref{tab:metricas}, $R^2 \approx 0{,}09$--$0{,}10$):
ao remover a variância de estrutura etária --- que estava correlacionada com
região/desenvolvimento e ``inflava'' o $R^2$ da taxa bruta ($\approx 0{,}23$) --- resta apenas
o sinal socioeconômico ``puro'', menor mas no sentido correto. O objetivo aqui é
\emph{inferência}, não predição.

\begin{table}[htbp]
\centering\small
\caption{Desempenho dos modelos (alvo: taxa padronizada; \emph{hold-out} por município).}
\label{tab:metricas}
\rowcolors{2}{cinzalin}{white}
\begin{tabular}{l r r r}
\rowcolor{cinzahdr}
\hdr{Modelo} & \hdr{$R^2$} & \hdr{RMSE} & \hdr{MAE} \\
\midrule
OLS (linear)   & 0,092 & 71,0 & 53,9 \\
Random Forest  & 0,091 & 71,0 & 54,1 \\
XGBoost        & 0,103 & 70,6 & 53,7 \\
\bottomrule
\end{tabular}
\end{table}

No painel completo (com controles de região e ano), o \textbf{Gini} permanece o achado central:
$+151$ ($p<0{,}001$) --- mais desigualdade, mais mortalidade. A renda \emph{per capita}
($-0{,}14$), a esperança de vida ($-4{,}2$) e a proporção de pobres ($-2{,}1$, colinear com
renda) também são significativas. O efeito do IDH-M é absorvido pelas \emph{dummies} de região
(N/NE são justamente as de menor IDH). Para isolá-lo, a Tabela~\ref{tab:coef} compara um modelo
transversal \emph{sem} controles de região, com alvo bruto vs.\ padronizado --- evidenciando a
correção do paradoxo.

\begin{table}[htbp]
\centering\small
\caption{Coeficiente do IDH-M em modelo transversal sem controles de região: a padronização
inverte o sinal para o esperado. \mbox{*** $p<0{,}001$;\ * $p<0{,}05$.}}
\label{tab:coef}
\rowcolors{2}{cinzalin}{white}
\begin{tabular}{l r r}
\rowcolor{cinzahdr}
\hdr{Variável} & \hdr{Alvo: taxa bruta} & \hdr{Alvo: taxa padronizada} \\
\midrule
IDH-M    & $-77{,}1$\,***  & $-126{,}5$\,*** \\
Gini     & $+43{,}6$\,*    & $+85{,}8$\,***  \\
Esp.\ de vida & $-0{,}1$   & $-3{,}5$\,***   \\
\bottomrule
\end{tabular}
\end{table}

Com a taxa padronizada, o IDH-M passa a ter efeito \textbf{negativo e fortemente significativo}
(mais desenvolvimento $\rightarrow$ menos mortalidade) e o Gini reforça --- exatamente os sinais
previstos pela hipótese.

\subsection{Clusters de municípios}

A clusterização produziu quatro perfis bem separados (ANOVA da taxa entre clusters altamente
significativa), renumerados por \textbf{IDH-M crescente} (cluster 0 = mais vulnerável)
--- Tabela~\ref{tab:clusters}. Reportamos as duas taxas: a coluna padronizada é a comparável.

\begin{table}[htbp]
\centering\small
\caption{Perfil dos quatro clusters de municípios (0 = menor IDH-M).}
\label{tab:clusters}
\rowcolors{2}{cinzalin}{white}
\begin{tabular}{c r r r c c c l}
\rowcolor{cinzahdr}
\hdr{Cl.} & \hdr{Mun.} & \hdr{Taxa bruta} & \hdr{Taxa padr.} & \hdr{IDH-M} & \hdr{Gini} & \hdr{Água\%} & \hdr{Região dom.} \\
\midrule
0 &  784 & 183 & 232 & 0,56 & 0,54 & 58 & NE/N \\
1 & 1.573 & 208 & 249 & 0,60 & 0,52 & 84 & NE \\
2 & 1.430 & 248 & 240 & 0,70 & 0,45 & 92 & SE/S \\
3 & 1.774 & 243 & 251 & 0,72 & 0,48 & 94 & SE/S \\
\bottomrule
\end{tabular}
\end{table}

Os clusters 0 e 1 reúnem municípios de menor IDH-M e pior saneamento, concentrados no
Nordeste e Norte; os clusters 2 e 3, os mais desenvolvidos do Sudeste e Sul. Pela taxa
\emph{bruta}, os desenvolvidos pareciam ter mortalidade maior (artefato etário); pela
\emph{padronizada}, as taxas se aproximam (232--251) e o cluster vulnerável 1 fica em patamar
comparável ao dos desenvolvidos --- a falsa vantagem das regiões pobres desaparece.

\subsection{Mineração: do que morre cada perfil}

A composição agregada de causas é \emph{semelhante} entre clusters --- todos dominados por
infarto, diabetes e neoplasias. O sinal socioeconômico aparece nas causas
\textbf{sobre-representadas} (Tabela~\ref{tab:lift}):

\begin{table}[htbp]
\centering\small
\caption{Causas mais sobre-representadas (\emph{lift}) por cluster.}
\label{tab:lift}
\rowcolors{2}{cinzalin}{white}
\begin{tabularx}{\textwidth}{l Y}
\rowcolor{cinzahdr}
\hdr{Cluster} & \hdr{Causas características (lift)} \\
\midrule
0 (vulnerável) & AVC (1,6$\times$), transtornos por álcool (1,6$\times$), diabetes, doença de Chagas, diarreia \\
1 & esquistossomose (2,5$\times$), doença de Chagas (1,7$\times$), álcool, infecções intestinais \\
2 (desenvolvido) & cardiopatia isquêmica crônica, infarto cerebral, neoplasia de cólon, hepatite crônica \\
3 & dengue, enfisema, hipertensão, neoplasias de esôfago e laringe \\
\bottomrule
\end{tabularx}
\end{table}

Há, portanto, uma \textbf{assinatura socioeconômica dupla}: perfis pobres ainda carregam
doenças infecciosas e negligenciadas (Chagas, esquistossomose, diarreia) que o
desenvolvimento deveria ter eliminado, enquanto perfis ricos concentram doenças crônicas. A
morte materna também segue o gradiente esperado, caindo de 0,65\% dos óbitos no cluster 0
para 0,31\% no cluster 3.

% ===================================================================
\section{Discussão e Limitações}

\textbf{Padronização por idade (resolvida).} A taxa \emph{bruta} produzia um resultado
paradoxal --- maior nos clusters desenvolvidos --- por confundimento etário: municípios
desenvolvidos têm população mais envelhecida na faixa 5--74 anos e acumulam mais mortes por
doenças crônicas. A padronização direta (com a estrutura etária do Censo 2022) corrige esse
viés: o gradiente regional se realinha (Nordeste passa a superar Sudeste/Sul) e o coeficiente
do IDH-M se torna negativo, como previsto. Permanecem duas ressalvas técnicas: (i) a
padronização \emph{direta} é sensível a ruído em municípios pequenos com poucas mortes por
faixa --- uma padronização \emph{indireta} (SMR) seria mais estável; (ii) o denominador por
idade é do Censo 2022, aplicado aos três anos (a estrutura etária varia pouco no período).

\textbf{Subnotificação em áreas remotas.} O cluster mais vulnerável (0, Norte/Nordeste rural)
apresenta a menor taxa padronizada, o que provavelmente reflete subnotificação de óbitos e
causas mal-definidas nessas áreas, e não menor mortalidade real --- uma limitação dos próprios
registros do SIM.

\textbf{Endogeneidade da despesa em saúde.} O coeficiente positivo da despesa \emph{per
capita} (mais gasto associado a mais mortes) não deve ser lido como causal: municípios com
maior carga de doença tendem a gastar mais (causalidade reversa). Isolar esse efeito exigiria
instrumentos ou defasagens temporais.

\textbf{Atlas DH defasado.} IDH-M, Gini e renda referem-se a 2010, enquanto os óbitos são de
2022--2024. Essa defasagem de medição atenua as associações e contribui para o $R^2$ modesto.

\textbf{Natureza ecológica.} Os clusters e as associações são de nível municipal. Não se pode
inferir, a partir deles, o comportamento de indivíduos (falácia ecológica).

% ===================================================================
\section{Conclusão}

O projeto entregou um pipeline reprodutível que integra cinco fontes públicas e sustenta três
análises complementares. As evidências apontam, de forma consistente, que a mortalidade por
causas evitáveis no Brasil tem natureza socioeconômica: a desigualdade (Gini) está
associada positivamente à taxa de forma robusta, e o \emph{tipo} de causa evitável varia
sistematicamente com o perfil do município --- doenças infecciosas e negligenciadas nos
territórios pobres, doenças crônicas nos desenvolvidos. A padronização por idade, incorporada
ao \emph{pipeline}, corrigiu o paradoxo das taxas brutas e alinhou os sinais da regressão à
hipótese (IDH-M negativo, Gini positivo). O aprofundamento futuro mais relevante é a
incorporação de variáveis da rede de atenção à saúde (leitos e médicos por habitante,
cobertura da atenção básica) e o tratamento da subnotificação em áreas remotas.

% ===================================================================
\section*{Referências}
\addcontentsline{toc}{section}{Referências}
\begin{itemize}[nosep]
  \item Malta DC et al. \emph{Lista de causas de mortes evitáveis por intervenções do SUS.}
        Epidemiol.\ Serv.\ Saúde, 16(4):233--244, 2007.
  \item Malta DC et al. \emph{Atualização da lista de causas de mortes evitáveis}
        [Nota técnica]. Epidemiol.\ Serv.\ Saúde, 19(2):173--176, 2010.
  \item DATASUS. \emph{Sistema de Informações sobre Mortalidade (SIM)} --- microdados CID-10.
  \item IBGE. \emph{Sidra} --- População e PIB municipal.
  \item Tesouro Nacional. \emph{SICONFI} --- Declaração de Contas Anuais (DCA).
  \item IPEA. \emph{Atlas do Desenvolvimento Humano} (Censo 2010).
\end{itemize}

\end{document}
